% -----------------------------------------------------------
\subsection{Gobernanza, Roles y Alcance Operativo}
% -----------------------------------------------------------

\subsubsection{Estructura Organizacional}

La plataforma soporta una estructura organizacional jerárquica compuesta por tres
niveles principales: cadena restaurantera, zona y sucursal. Esta organización permite
asignar permisos y consultar información de acuerdo con el alcance operativo de cada
usuario colaborador.

\begin{itemize}[leftmargin=1.5em]
  \item \textbf{Cadena restaurantera:} entidad raíz que agrupa zonas y sucursales.
  \item \textbf{Zona:} agrupación geográfica u operativa de sucursales dentro de una
  cadena.
  \item \textbf{Sucursal:} unidad operativa donde se gestionan mesas, órdenes, cocina,
  liquidación interna y datos analíticos.
\end{itemize}

La asignación de permisos se realiza mediante roles contextuales. Cada usuario colaborador
opera a nivel plataforma, cadena, zona o restaurante, según el valor de
\texttt{UserRole.contextType} y las llaves contextuales asociadas.

\subsubsection{Roles de Usuarios}

El sistema distingue entre usuarios colaboradores y actores temporales de consumo. Los
usuarios colaboradores pertenecen a la operación del restaurante y se autentican mediante
cuenta. Los actores temporales, como Anfitrión y Comensal, acceden mediante la sesión de
mesa virtual y no forman parte del RBAC administrativo.

\begin{table}[H]
\caption{Roles y alcance operativo}
\label{tab:roles_alcance}
\centering
\small
\begin{tabular}{p{4cm} p{10cm}}
\hline
\thead{Rol} & \thead{Alcance} \\
\hline
Admin Cadena & Administración de la cadena, zonas, sucursales, usuarios, menú global
y dashboard consolidado. \\
Gerente Zona & Gestión y consulta de las sucursales pertenecientes a su zona. \\
Gerente Sucursal & Gestión operativa de una sucursal específica: menú local, mesas,
personal, órdenes y liquidaciones internas. \\
Mesero & Gestión de mesas, sesiones, órdenes y cierre operativo dentro de su sucursal. \\
Cocina/Chef & Visualización y actualización de órdenes asignadas a estaciones de
preparación. \\
Anfitrión & Activación temporal de la mesa virtual y administración básica de acceso
de comensales. \\
Comensal & Consulta de menú, generación de órdenes y registro de aportaciones internas
dentro de una mesa virtual. \\
\hline
\end{tabular}
\end{table}

Además del control de permisos en la capa de aplicación, el diseño contempla políticas
de seguridad a nivel de fila (RLS) en la base de datos para restringir el acceso según
el alcance organizacional del usuario. De esta forma, los permisos no dependen únicamente
de la interfaz, sino también de reglas de aislamiento en la capa de persistencia.

\subsubsection{Reglas de Asignación de Roles a Niveles}

\begin{itemize}[leftmargin=1.5em]
  \item El usuario \textbf{Admin Cadena} opera a nivel cadena.
  \item El \textbf{Gerente Zona} opera sobre las sucursales pertenecientes a su zona.
  \item El \textbf{Gerente Sucursal}, \textbf{Mesero} y \textbf{Cocina/Chef} operan a
  nivel sucursal.
  \item El \textbf{Anfitrión} y el \textbf{Comensal} existen únicamente dentro del
  contexto temporal de una \texttt{DiningSession}.
\end{itemize}

La Tabla~\ref{tab:permisos} presenta la matriz de permisos por rol del sistema.

\setstretch{1.2}
\setlength{\tabcolsep}{4pt}
\begin{longtable}{>{\raggedright\arraybackslash}p{10cm}
                >{\centering\arraybackslash}p{0.8cm}
                >{\centering\arraybackslash}p{1.1cm}
                >{\centering\arraybackslash}p{1.1cm}
                >{\centering\arraybackslash}p{1cm}
                >{\centering\arraybackslash}p{1.1cm}
                >{\centering\arraybackslash}p{0.8cm}}
\caption{Matriz de Permisos por Rol} \label{tab:permisos} \\
\hline
\thead{Acción / Permiso} &
\thead{Chef} &
\thead{Comensal} &
\thead{Anfi-trión} &
\thead{Mesero} &
\thead{Gerente} &
\thead{Admin} \\
\hline
\endfirsthead
\caption[]{TABLA \thetable{} (Continuación)} \\
\hline
\thead{Acción / Permiso} &
\thead{Chef} &
\thead{Comensal} &
\thead{Anfi-trión} &
\thead{Mesero} &
\thead{Gerente} &
\thead{Admin} \\
\hline
\endhead
\hline
\endfoot
Dar de alta Zonas                                      &   &   &   &   &   & X \\ \rowcolor{altrow}
Dar de alta usuarios administrativos de cadena         &   &   &   &   &   & X \\
Dar de alta usuarios con otros roles                   &   &   &   &   & X & X \\ \rowcolor{altrow}
Acceso al dashboard analítico                          &   &   &   &   & X & X \\
Dar de alta Sucursales                                 &   &   &   &   & X & X \\ \rowcolor{altrow}
Agregar, eliminar o editar menús y productos           &   &   &   &   & X &   \\
Crear, editar y liberar mesas físicas                  &   &   &   &   & X & X \\ \rowcolor{altrow}
Visualizar mesas físicas y características             &   &   &   & X & X & X \\
Inicializar mesa virtual                               &   &   &   & X & X & X \\ \rowcolor{altrow}
Activar mesa virtual y crear contraseña de acceso      &   &   & X &   &   &   \\
Compartir acceso temporal a la mesa virtual            &   & X & X &   &   &   \\ \rowcolor{altrow}
Generar órdenes individuales                           &   & X & X &   &   &   \\
Generar órdenes compartidas                            &   & X & X &   &   &   \\ \rowcolor{altrow}
Visualizar órdenes y estado de cocina                  & X &   &   & X & X & X \\
Visualizar desglose de consumo por comensal            &   & X & X & X & X & X \\ \rowcolor{altrow}
Registrar aportación interna o liquidar consumo        &   & X & X &   &   &   \\
Visualizar estado de liquidación interna               &   & X & X & X & X & X \\
\hline
\end{longtable}
\setstretch{1.5}

\noindent\textit{Nota 1: La columna Gerente agrupa a Gerente Zona y Gerente Sucursal.
Los permisos se aplican respetando el alcance definido por \texttt{UserRole.contextType};
por ejemplo, un Gerente Zona opera solo sobre las sucursales de su zona y un Gerente
Sucursal solo sobre su sucursal.}

\noindent\textit{Nota 2: La creación de nuevos usuarios administrativos de cadena debe
quedar registrada en \texttt{AuditLog} y puede restringirse por política interna.}

\noindent\textit{Nota 3: La contraseña de mesa se almacena como hash; el sistema no debe
mostrarla nuevamente después de su creación. El personal de cocina visualiza y actualiza
ítems hasta el estado \texttt{LISTA}. El mesero confirma la entrega al comensal mediante
el cambio de \texttt{LISTA} a \texttt{ENTREGADA}.}

\medskip

La Tabla~\ref{tab:alta_usuarios} define qué roles puede crear cada tipo de usuario
colaborador del sistema.

\begin{table}[H]
\caption{Reglas de creación de usuarios colaboradores}
\label{tab:alta_usuarios}
\centering
\small
\begin{tabular}{p{4cm} p{10cm}}
\hline
\thead{Rol creador} & \thead{Roles que puede crear} \\
\hline
Admin Cadena & Gerente Zona, Gerente Sucursal, Mesero y Cocina/Chef dentro de cualquier
sucursal de la cadena. \\
Gerente Zona & Gerente Sucursal, Mesero y Cocina/Chef dentro de las sucursales de
su zona. \\
Gerente Sucursal & Mesero y Cocina/Chef dentro de su sucursal. \\
Mesero & No puede crear usuarios. \\
Cocina/Chef & No puede crear usuarios. \\
\hline
\end{tabular}
\end{table}

% -----------------------------------------------------------
\subsection{Delimitación Técnica del Diseño}
% -----------------------------------------------------------

\subsubsection{Dentro del Alcance}

\begin{itemize}[leftmargin=1.5em]
  \item Alta y gestión de mesas físicas para una sola área o piso por sucursal.
  \item Gestión de mesas virtuales con activación por QR firmado, contraseña de acceso,
  órdenes y liquidación interna.
  \item Módulo de órdenes individuales y compartidas para comensales.
  \item Tablero KDS en tiempo real para cocina.
  \item Dashboard analítico para gerentes y administrador de cadena.
  \item Gestión de menú, categorías, productos, variantes y modificadores.
  \item Gestión de usuarios colaboradores con jerarquía de permisos.
  \item Módulo de liquidación interna con control de concurrencia y constancia interna
  de estado de cuenta.
  \item Pipeline analítico batch con capas Bronze, Silver y Gold.
  \item Tablas Gold para métricas de consumo, productos, tiempos de servicio y
  estimaciones exploratorias de demanda.
\end{itemize}

\subsubsection{Fuera del Alcance}

\begin{itemize}[leftmargin=1.5em]
  \item Seccionamiento de mesas físicas por múltiples áreas o pisos.
  \item Recuperación autónoma de contraseña de mesa virtual por parte del comensal sin
  intervención del personal autorizado.
  \item Flujo de recuperación de contraseña para usuarios colaboradores.
  \item Integración con pasarelas de pago bancarias.
  \item Emisión de CFDI o facturación fiscal.
  \item Descarga avanzada de reportes XLSX o PDF de consumo, operación y métricas
  analíticas.
  \item Clúster distribuido dedicado para procesamiento analítico de alto volumen.
  \item Modelos predictivos avanzados de \textit{machine learning} o recomendadores
  personalizados.
\end{itemize}

\subsubsection{Consideraciones de Diseño}

\begin{itemize}[leftmargin=1.5em]
  \item Las aportaciones internas concurrentes deben serializarse para evitar
  sobreliquidación.
  \item La sesión de mesa virtual debe validar QR firmado, contraseña y estado activo.
  \item El cierre operativo de mesa debe liberar mesas físicas y cerrar la sesión lógica.
  \item Las métricas históricas deben consultarse desde tablas Gold para evitar saturar
  la base OLTP.
  \item Toda acción crítica debe registrarse en \texttt{AuditLog} para fortalecer la
  trazabilidad operativa y facilitar auditorías internas.
\end{itemize}

\subsubsection{Trabajo Futuro}

\begin{itemize}[leftmargin=1.5em]
  \item Integración con pasarelas de pago certificadas.
  \item Descarga de reportes XLSX/PDF por período de tiempo.
  \item Flujo avanzado de recuperación de acceso para anfitrión, validado por mesero
  o gerente.
  \item Soporte multiárea o multipiso para planos de restaurante.
  \item Modelo de \textit{machine learning} para recomendación de platillos.
  \item Botón para llamar al mesero desde la aplicación del comensal.
  \item Notificaciones \textit{push} al mesero para acciones requeridas.
\end{itemize}
